\documentclass[12pt]{article}
\usepackage[utf8]{inputenc}
\usepackage[T1]{fontenc}
\usepackage{hyperref}
\title{Jemez Language Resources for K-12 Curriculum Development}
\author{Compiled by Claude Sonnet 5 (claude-sonnet-5) for Andr\'as Kornai}
\date{}
\begin{document}
\maketitle

\section{Introduction}
Jemez (Towa; ISO 639-3 tow) is a Kiowa-Tanoan language spoken at Jemez Pueblo (Walatowa), New Mexico. Speaker counts vary by source (Wikipedia gives 1,800 in one version and 3,000 in another). The search found no published modern dictionary, textbook or primer. According to the sources found, the community does not permit the language to be written down, so any curriculum work should go through the Pueblo of Jemez.

\section{Primary Educational Resources}
\begin{itemize}
\item Walatowa Head Start Language Immersion Program. Reported to have started a Towa immersion programme in 2007. News coverage says immersion pilots now extend up to eighth grade, funded by public and private grants. \url{https://www.santafenewmexican.com/news/local_news/indigenous-language-immersion-schools-aim-to-reverse-cultural-loss/article_cb68e890-f458-11ef-8072-e7099c694d6a.html}
\item Yahoo News copy of the same report on Indigenous-language immersion schools. \url{https://www.yahoo.com/news/indigenous-language-immersion-schools-aim-043500014.html}
\item Presentation by Kevin Shendo on education systems and programs at Jemez Pueblo (First Nations Development Institute webinar). \url{https://www.firstnations.org/wp-content/uploads/2024/12/NAHM-Webinar-3-Education-Systems-and-Programs-at-Jemez-Pueblo-with-Kevin-Shendo.pdf}
\end{itemize}

\section{Linguistic and Reference Materials}
\begin{itemize}
\item Smithsonian National Anthropological Archives, MS 1657: a 1907 vocabulary of about 780 Jemez words with English equivalents collected by Albert B. Reagan, with a grammar treatise. It is an archival manuscript, not a published dictionary. \url{https://sirismm.si.edu/EADpdfs/NAA.MS1657.pdf}
\item Wikipedia overview of the Jemez language. \url{https://en.wikipedia.org/wiki/Jemez_language}
\end{itemize}

\section{Religious and Community Resources}
No scripture or religious materials were found in the search results. Language and religion are closely tied at the Pueblo, and access to such material is a matter for the community.
\begin{itemize}
\item Char-Koosta News article on Jemez Pueblo's efforts to revitalize language and culture. \url{https://www.charkoosta.com/news/jemez-pueblo-strives-to-revitalize-language-culture/article_0ce0fd5e-a8cd-11e9-8c34-932b8f4d412c.html}
\end{itemize}

\section{Digital and Audio Resources}
No audio, video or software resources were found in the search results.

\section{Curriculum Development Context}
Immersion programmes exist at Jemez Pueblo, which contradicts the survey note that the language is not known to be taught in schools. Published teaching materials were not found. The Pueblo of Jemez Department of Education is the natural partner. The sources say there is no common written form and that tribal rules do not allow transcription. Any written curriculum therefore depends on community permission.
\end{document}